% ============================================================================
% Chapitre 3 — Analyse et spécification des besoins
% ============================================================================
\chapter{Analyse et spécification des besoins}

\section*{Introduction}

Ce chapitre formalise les exigences du système \caimp{} en besoins fonctionnels
et non fonctionnels, identifie les acteurs, et présente les modèles UML qui
guident la phase de conception. La démarche d'analyse s'est appuyée sur l'étude
comparative du chapitre précédent, complétée par des scénarios d'usage concrets.

\section{Identification des acteurs}

Le système \caimp{} implique quatre catégories d'acteurs :

\begin{description}[leftmargin=2.5cm, style=nextline]
  \item[\textbf{Administrateur}] Rôle le plus privilégié. Gère les organisations,
    les utilisateurs, les serveurs, et configure la plateforme. Accède à l'\api{}
    d'administration (port 8001).
  \item[\textbf{Opérateur}] Utilisateur technique qui consulte les tableaux de
    bord, gère les runbooks, configure les règles d'alerte, et interprète les
    analyses \textsc{ai}. Rôle \texttt{devops} en base.
  \item[\textbf{Viewer}] Lecture seule. Consulte les dashboards et les rapports
    PDF. Rôle conçu pour les parties prenantes non techniques (managers, clients).
  \item[\textbf{Agent système}] Acteur non humain. Le processus démon
    \texttt{caimp-agent} s'exécute sur les serveurs surveillés. Il s'authentifie
    par certificat mTLS et envoie les métriques via OTLP.
\end{description}

\section{Besoins fonctionnels}

\begin{enumerate}[label=\textbf{BF\arabic*}, leftmargin=2.2cm, itemsep=4pt]
  \item L'agent doit collecter les métriques CPU, mémoire et disque toutes les
    15 secondes et les transmettre via OTLP/HTTP.
  \item L'agent doit s'auto-inscrire auprès de la plateforme à l'aide d'un jeton
    d'enrôlement à usage unique.
  \item L'agent doit utiliser un certificat mTLS délivré par l'autorité de
    certification interne pour authentifier chaque connexion.
  \item L'agent doit détecter et reporter automatiquement les événements Docker
    (démarrage, arrêt, tirage d'image).
  \item L'agent doit surveiller les journaux d'installation de paquets (dpkg, yum)
    et les événements SSH depuis les journaux d'authentification.
  \item La plateforme doit détecter les anomalies en temps réel à l'aide d'au moins
    trois algorithmes distincts (seuil, Z-score, taux de variation).
  \item Les anomalies détectées doivent être publiées sur NATS JetStream pour
    traitement asynchrone par les services consommateurs.
  \item Le moteur d'analyse IA doit interroger la base vectorielle pour retrouver
    des incidents similaires passés et les runbooks pertinents.
  \item Le moteur d'analyse IA doit soumettre un prompt structuré à Ollama et
    valider la réponse JSON avant persistance.
  \item L'explication IA doit contenir : l'explication en langage naturel, la
    cause racine probable, l'action recommandée, et le niveau de confiance.
  \item Le moteur de corrélation doit interroger la table \texttt{change\_events}
    dans la fenêtre de 30 minutes précédant chaque anomalie.
  \item Le moteur de corrélation doit scorer chaque changement par proximité
    temporelle et type, et produire un classement des causes probables.
  \item Le moteur de prévision doit calculer les prévisions Holt-Winters sur
    24 heures pour chaque combinaison serveur/métrique.
  \item Le moteur de prévision doit calculer le \textit{time to critical} ---
    le délai prévu avant dépassement de seuil à 90~\%.
  \item La déduplication d'alertes doit supprimer les alertes identiques
    (même serveur, même métrique, même détecteur) répétées dans les 15 minutes.
  \item L'\api{} de requête doit exposer un endpoint \texttt{/dashboard/summary}
    agrégant le statut global, les problèmes potentiels et les actions recommandées.
  \item L'interface utilisateur doit présenter en première page le statut global,
    les problèmes potentiels triés par criticité, et les actions recommandées
    numérotées.
  \item La plateforme doit permettre la gestion des runbooks (créer, modifier,
    supprimer) avec ré-indexation automatique des embeddings.
  \item Le module de rapports doit générer un rapport PDF multi-sections lisible
    par des non-techniciens, avec logo de l'organisation.
  \item L'interface doit proposer un chatbot SSE permettant de poser des questions
    en langage naturel sur les incidents et les serveurs.
  \item La plateforme doit recevoir les webhooks GitHub et GitLab pour les
    événements de déploiement.
  \item L'administrateur doit pouvoir configurer les canaux de notification
    (Slack, PagerDuty) pour les alertes critiques.
  \item Les mises à jour d'état doivent être poussées en temps réel vers les
    navigateurs connectés via WebSocket.
  \item La plateforme doit supporter l'isolation multi-tenant : les données d'une
    organisation ne doivent jamais être accessibles par une autre.
  \item La plateforme doit fournir un endpoint de santé (\texttt{/health}) pour
    chaque service, utilisé par les healthchecks Docker.
\end{enumerate}

\section{Besoins non fonctionnels}

\subsection{Performance}

\begin{itemize}
  \item \textbf{Latence d'ingestion :} Le Telemetry Writer doit traiter un batch
    de métriques en moins de 50~ms en conditions normales.
  \item \textbf{Throughput :} La plateforme doit supporter au moins 1~000 points
    de métriques par seconde sur un déploiement standard.
  \item \textbf{Temps de réponse API :} Les endpoints de lecture (GET) doivent
    répondre en moins de 300~ms pour un volume de 50 serveurs.
  \item \textbf{Latence IA :} Une explication complète doit être générée en moins
    de 30 secondes (dépendant du modèle LLM et du matériel).
\end{itemize}

\subsection{Sécurité}

\begin{itemize}
  \item Authentification de chaque agent par certificat mTLS (TLS 1.3) ;
  \item Tokens d'enrôlement à usage unique, expirés après 30 minutes ;
  \item Jetons JWT HS256 avec expiration courte (15 minutes) pour les \api{} ;
  \item Isolation des données par organisation via RLS PostgreSQL ;
  \item Chiffrement AES-256-GCM des secrets Splunk en base de données ;
  \item Journalisation immuable de toutes les actions administratives (\texttt{audit\_log}).
\end{itemize}

\subsection{Disponibilité et résilience}

\begin{itemize}
  \item Objectif de disponibilité : 99,5~\% (temps d'arrêt maximum de 3h38/mois) ;
  \item Healthchecks Docker sur chaque service avec politique de redémarrage
    \texttt{unless-stopped} ;
  \item Buffering local de l'agent sur SQLite en cas d'indisponibilité temporaire
    de la plateforme ;
  \item Messages NATS JetStream persistés sur disque (durabilité garantie).
\end{itemize}

\subsection{Maintenabilité}

\begin{itemize}
  \item Chaque service est autonome, déployable et testable indépendamment ;
  \item Migrations de base de données versionnées avec \texttt{golang-migrate} ;
  \item Conventions de commit standardisées (\textit{Conventional Commits}) ;
  \item Documentation inline pour les fonctions publiques.
\end{itemize}

\section{Diagramme de cas d'utilisation global}

\begin{figure}[H]
\centering
\begin{tikzpicture}[
  actor/.style={draw=NavyBlue, fill=NavyBlue!10, rounded corners=3pt,
    minimum width=1.8cm, minimum height=0.7cm, font=\small\bfseries},
  uc/.style={draw=TealBlue, fill=TealBlue!8, ellipse,
    minimum width=3.4cm, minimum height=0.9cm, font=\small, align=center},
  sysbox/.style={draw=NavyBlue!50, dashed, rounded corners=6pt, fill=NavyBlue!3}
]

% Système CAIMP
\node[sysbox, minimum width=13cm, minimum height=11cm] (sys) at (5.5,0) {};
\node[above, font=\small\bfseries\color{NavyBlue}] at (sys.north) {Système CAIMP};

% Acteurs gauche
\node[actor] (admin)  at (-1.2, 3.0)   {Admin};
\node[actor] (ops)    at (-1.2, 0.5)   {Opérateur};
\node[actor] (viewer) at (-1.2,-2.0)   {Viewer};
\node[actor] (agent)  at (12.5, 0.5)   {Agent};

% Cas d'utilisation
\node[uc] (login)     at (3.0, 4.5)    {S'authentifier};
\node[uc] (usrmgr)    at (7.0, 4.5)    {Gérer utilisateurs};
\node[uc] (srvmgr)    at (7.0, 3.0)    {Gérer serveurs};
\node[uc] (dash)      at (3.0, 1.8)    {Consulter\\dashboard};
\node[uc] (alerts)    at (7.0, 1.8)    {Voir alertes\\en temps réel};
\node[uc] (runbooks)  at (3.0, 0.2)    {Gérer runbooks};
\node[uc] (chat)      at (7.0, 0.2)    {Utiliser le chat IA};
\node[uc] (report)    at (3.0,-1.4)    {Exporter rapport\\PDF};
\node[uc] (forecast)  at (7.0,-1.4)    {Voir prévisions};
\node[uc] (metrics)   at (5.0,-3.0)    {Envoyer métriques\\(OTLP)};
\node[uc] (enroll)    at (9.0,-3.0)    {S'enrôler (mTLS)};

% Flèches acteurs → UC
\draw[->, NavyBlue] (admin)  -- (login);
\draw[->, NavyBlue] (admin)  -- (usrmgr);
\draw[->, NavyBlue] (admin)  -- (srvmgr);
\draw[->, NavyBlue] (ops)    -- (login);
\draw[->, NavyBlue] (ops)    -- (dash);
\draw[->, NavyBlue] (ops)    -- (alerts);
\draw[->, NavyBlue] (ops)    -- (runbooks);
\draw[->, NavyBlue] (ops)    -- (chat);
\draw[->, NavyBlue] (ops)    -- (forecast);
\draw[->, NavyBlue] (viewer) -- (login);
\draw[->, NavyBlue] (viewer) -- (dash);
\draw[->, NavyBlue] (viewer) -- (report);
\draw[->, NavyBlue] (agent)  -- (metrics);
\draw[->, NavyBlue] (agent)  -- (enroll);

\end{tikzpicture}
\caption{Diagramme de cas d'utilisation global de \caimp{}}
\label{fig:ucd-global}
\end{figure}

\section{Diagrammes de séquence}

\subsection{Enrôlement d'un agent}

\begin{figure}[H]
\centering
\begin{tikzpicture}[
  lifeline/.style={draw=NavyBlue!50, dashed, thick},
  actor/.style={draw=NavyBlue, fill=NavyBlue!15, minimum width=2.2cm,
    minimum height=0.6cm, font=\footnotesize\bfseries, align=center},
  msg/.style={->, draw=TealBlue, thick},
  retmsg/.style={->, draw=DarkGray, dashed},
  note/.style={draw=AccentOrange!60, fill=AccentOrange!10,
    font=\tiny, align=center, rounded corners=3pt}
]
\def\dt{0.9}

% Acteurs
\node[actor] (admin) at (0,0)    {Admin\\(Browser)};
\node[actor] (adm)   at (3.5,0)  {Admin API\\:8001};
\node[actor] (db)    at (7.0,0)  {PostgreSQL};
\node[actor] (agt)   at (10.5,0) {Agent\\(cible)};

% Lignes de vie
\draw[lifeline] (admin.south) -- ++(0,-10.5);
\draw[lifeline] (adm.south)   -- ++(0,-10.5);
\draw[lifeline] (db.south)    -- ++(0,-10.5);
\draw[lifeline] (agt.south)   -- ++(0,-10.5);

% Séquence
\draw[msg] ($(admin.south)+(0,-0.5)$) -- node[above,font=\tiny]{POST /servers/enroll}
  ($(adm.south)+(0,-0.5)$);
\draw[msg] ($(adm.south)+(0,-1.2)$) -- node[above,font=\tiny]{INSERT enrollment\_token}
  ($(db.south)+(0,-1.2)$);
\draw[retmsg] ($(db.south)+(0,-2.0)$) -- node[above,font=\tiny]{token, server\_id}
  ($(adm.south)+(0,-2.0)$);
\draw[retmsg] ($(adm.south)+(0,-2.8)$) -- node[above,font=\tiny]{\{enrollment\_token\}}
  ($(admin.south)+(0,-2.8)$);

\node[font=\small\bfseries\color{TealBlue}] at (5.5,-3.5) {— L'admin copie le token sur le serveur cible —};

\draw[msg] ($(agt.south)+(0,-4.2)$) -- node[below,font=\tiny]{python main.py --enroll TOKEN}
  ++(0,0);
\draw[msg] ($(agt.south)+(0,-4.8)$) -- node[above,font=\tiny]{POST /servers/\{id\}/enroll + token}
  ($(adm.south)+(0,-4.8)$) node[below left, note]{Vérifie token\\+ expiry};
\draw[msg] ($(adm.south)+(0,-5.8)$) -- node[above,font=\tiny]{Génère cert mTLS (ECDSA P-256)}
  ($(db.south)+(0,-5.8)$);
\draw[retmsg] ($(db.south)+(0,-6.6)$) -- node[above,font=\tiny]{cert\_pem, ca\_cert}
  ($(adm.south)+(0,-6.6)$);
\draw[retmsg] ($(adm.south)+(0,-7.4)$) -- node[above,font=\tiny]{\{cert, key, ca\_cert, server\_id\}}
  ($(agt.south)+(0,-7.4)$);
\draw[msg] ($(agt.south)+(0,-8.2)$) -- node[above,font=\tiny]{Sauvegarde cert local, démarre OTLP}
  ++(0,0);

\node[font=\footnotesize\bfseries\color{SuccessGreen}] at (5.5,-9.0)
  {Agent enrôlé — métriques commencent à arriver};
\end{tikzpicture}
\caption{Séquence d'enrôlement d'un agent de supervision}
\label{fig:seq-enrollment}
\end{figure}

\subsection{Détection d'anomalie et explication IA}

\begin{figure}[H]
\centering
\begin{tikzpicture}[
  lifeline/.style={draw=NavyBlue!50, dashed, thick},
  actor/.style={draw=NavyBlue, fill=NavyBlue!15, minimum width=2cm,
    minimum height=0.6cm, font=\tiny\bfseries, align=center},
  msg/.style={->, draw=TealBlue, thick, font=\tiny},
  retmsg/.style={->, draw=DarkGray, dashed, font=\tiny}
]
\node[actor] (agt)  at (0,0)    {Agent};
\node[actor] (tw)   at (2.5,0)  {Telemetry\\Writer};
\node[actor] (nats) at (5.0,0)  {NATS};
\node[actor] (ai)   at (7.5,0)  {AI Worker};
\node[actor] (oll)  at (10.0,0) {Ollama};
\node[actor] (ws)   at (12.5,0) {WS Gwy};

\foreach \x in {agt,tw,nats,ai,oll,ws}
  \draw[lifeline] (\x.south) -- ++(0,-9.5);

\draw[msg] ($(agt.south)+(0,-0.6)$) -- node[above]{OTLP batch métriques}
  ($(tw.south)+(0,-0.6)$);
\draw[msg] ($(tw.south)+(0,-1.3)$) -- node[above]{INSERT metrics (batch)}
  ($(nats.south)+(0,-1.3)$) node[right,font=\tiny,color=DarkGray]{TimescaleDB};
\draw[msg] ($(tw.south)+(0,-2.1)$) -- node[above]{Detect() → anomaly}
  ($(nats.south)+(0,-2.1)$);
\draw[msg] ($(nats.south)+(0,-2.9)$) -- node[above]{anomaly.detected.*}
  ($(ai.south)+(0,-2.9)$);
\draw[msg] ($(ai.south)+(0,-3.7)$) -- node[above]{SELECT history + RAG}
  ($(oll.south)+(0,-3.7)$) node[right,font=\tiny,color=DarkGray]{PostgreSQL};
\draw[msg] ($(ai.south)+(0,-4.5)$) -- node[above]{POST /api/generate (prompt)}
  ($(oll.south)+(0,-4.5)$);
\draw[retmsg] ($(oll.south)+(0,-5.3)$) -- node[above]{\{explication, cause, action\}}
  ($(ai.south)+(0,-5.3)$);
\draw[msg] ($(ai.south)+(0,-6.1)$) -- node[above]{INSERT ai\_explanations}
  ($(nats.south)+(0,-6.1)$) node[right,font=\tiny,color=DarkGray]{PostgreSQL};
\draw[msg] ($(ai.south)+(0,-6.9)$) -- node[above]{ai.explanation.ready.*}
  ($(nats.south)+(0,-6.9)$);
\draw[msg] ($(nats.south)+(0,-7.7)$) -- node[above]{fan-out}
  ($(ws.south)+(0,-7.7)$);
\draw[msg] ($(ws.south)+(0,-8.5)$) -- node[above]{WebSocket push → Browser}
  ++(0,0);
\end{tikzpicture}
\caption{Séquence de détection d'anomalie et génération d'explication IA}
\label{fig:seq-anomaly}
\end{figure}

\section{Matrice de traçabilité des besoins}

\begin{table}[H]
\centering
\caption{Extrait de la matrice de traçabilité besoins fonctionnels -- composants}
\label{tab:tracabilite}
\small
\begin{tabularx}{\textwidth}{lXll}
\toprule
\textbf{BF} & \textbf{Besoin} & \textbf{Composant} & \textbf{Module} \\
\midrule
BF1  & Collecte métriques OTLP        & \servicename{caimp-agent}      & \texttt{collectors.py} \\
BF2  & Enrôlement par jeton            & \servicename{admin-api}        & \texttt{enrollment.py} \\
BF3  & mTLS agent                      & \servicename{otelcol / tw}     & Config mTLS \\
BF6  & Détection anomalies (3 algos)   & \servicename{telemetry-writer} & \texttt{detector.go} \\
BF7  & Publication NATS                & \servicename{telemetry-writer} & \texttt{publisher.go} \\
BF8--BF10 & Explication IA            & \servicename{ai-worker}        & \texttt{prompt\_builder.py} \\
BF11--BF12 & Corrélation changements  & \servicename{correlation-engine} & \texttt{main.py} \\
BF13--BF14 & Prévision Holt-Winters   & \servicename{forecast-engine}  & \texttt{forecaster.py} \\
BF15 & Déduplication alertes           & \servicename{alert-engine}     & \texttt{worker.py} \\
BF16--BF17 & Dashboard answers-first  & \servicename{query-api}        & \texttt{dashboard.py} \\
BF18 & Gestion runbooks + RAG          & \servicename{query-api}        & \texttt{incidents.py} \\
BF19 & Rapport PDF                     & \servicename{query-api}        & \texttt{reports.py} \\
BF21 & Webhooks GitHub/GitLab          & \servicename{change-tracker}   & \texttt{main.py} \\
BF23 & WebSocket push                  & \servicename{ws-gateway}       & \texttt{main.py} \\
BF24 & Isolation multi-tenant          & PostgreSQL                     & Migrations RLS \\
\bottomrule
\end{tabularx}
\end{table}

\section*{Conclusion}

L'analyse des besoins a permis d'identifier 25 besoins fonctionnels couvrant
la totalité du périmètre de \caimp{}, et de formaliser les exigences non
fonctionnelles en matière de performance, sécurité, disponibilité et
maintenabilité. Les diagrammes UML présentés constituent le contrat entre
l'analyse et la conception qui fait l'objet du chapitre suivant.
